\textbf{Non-triviality.} A free massive theory also has a gap. My first criterion (a nonzero connected four-point function of
$\mathrm{tr}\,F^2$) holds in free theories too, and I withdrew it~\cite{ClosedSet}. The fingerprint I now seek is logarithmic
running: by H6, $|x|^8\langle\Theta(x)\Theta(0)\rangle\sim C_O(N^2-1)/\bigl(16\ln^2(1/|x|\Lambda)\bigr)$ as
$|x|\to0$, an inverse square logarithm, unlike a nonzero dimension-four Wick polynomial in finitely many canonical
free fields with a nonvanishing leading scaling term. A \emph{generalized free field} (Gaussian with an arbitrary two-point function) could still mimic it; a
nonzero connected three-point function excludes that, and the leading one is $(b_0\bar g^2/32\pi^2)^3$ times the free
Wick triangle $T_O$ of three canonically normalized $O$ insertions, a finite computation I have not done.

\begin{claim}{CONDITIONAL}{Two tests of non-triviality}\label{b:nontrivial}
If H6 holds and $T_O\not\equiv0$, these correlations exclude the stated finite canonical Wick-polynomial comparison
and exclude $\Theta$ itself being a generalized free field. They do not alone classify the full observable algebra:
a non-Gaussian composite of Gaussian fields can have a nonzero connected three-point function.
Non-triviality of the reconstructed theory remains to be established with H6, not by that three-point test alone.
\evidence{The $1/16$ coefficient follows from one-loop running. The scope of the free-field comparison matters;
the leading triangle is still uncomputed~\cite{ClosedSet}.}
\lines{26, 25}
\end{claim}
